\section{Introduction}
\label{sec:introduction}

In 1988, David Godfrey examined Voyager images of Saturn's north pole and
found a hexagon \citep{Godfrey1988}.  Not an approximate hexagon, not a
roughly six-sided shape, but a persistent, regular, six-sided polygon
inscribed in the jet stream at $76^\circ$\,N---stable across decades,
seasons, and the passage of the Cassini mission through 2014.  Thirty
years later, the Juno spacecraft revealed that Jupiter's poles harbor
their own polygonal arrangements: eight cyclones in an octagonal ring at
the north pole, five in a pentagonal ring at the south
\citep{Adriani2018}.  Laboratory rotating tanks produce similar patterns
\citep{BarbosaAguiar2010}.

These observations raise a question that is simple to state and
surprisingly difficult to answer: \emph{why do rotating fluids make
polygons?}

The question has two parts.  The first---\emph{how} does a specific
polygon form?---has received considerable attention.  Numerical
simulations reproduce Saturn's hexagon from jet instabilities
\citep{Morales2011}; vortex dynamics models capture Jupiter's cyclone
rings; Rossby wave theory explains the wavenumber selection $n = 6$.
The second part---\emph{why} should different physical mechanisms on
different planets all produce the same geometric outcome?---has not
been addressed.

This paper answers both parts.  The answer has five steps:
\begin{enumerate}
  \item \textbf{Geometry.}  The M\"obius conjugacy theorem places
    \emph{circular} orbits at the locus $|\lam| = 1$ in the complex
    multiplicative group $\Cstar$, where logarithmic spirals transition
    to circles.  A separate geometric step (Proposition~\ref{prop:polygon})
    shows that imposing $\mathbb{Z}_N$ azimuthal symmetry converts these
    circles into $N$-gon jet meanders.

  \item \textbf{Branch disconnection.}  Within the WKB approximation,
    the Rayleigh--Kuo eigenvalue equation has two branches---elliptic
    ($\sigma = 0$, polygonal) and hyperbolic ($\alpha = 0$, radial)---that
    are disconnected.  Smooth perturbations of the jet speed cannot move
    the system between branches at leading order.

  \item \textbf{Energy maximum.}  Among $N$ point vortices on a ring, the
    regular $N$-gon maximizes the interaction energy: $H''(0) < 0$ for
    all $N \ge 3$ (proven via a closed-form formula that is manifestly
    negative).

  \item \textbf{Onsager condensation.}  In Onsager's (1949) statistical
    mechanics of two-dimensional vortices, systems at negative temperature
    are driven toward energy maxima.  The polygon is therefore the most
    probable macrostate in the large-scale limit of 2D turbulence.

  \item \textbf{Atmospheric thermostat.}  Small-scale turbulence in
    planetary atmospheres provides the effective mixing that enables
    Onsager relaxation.  The polygon is the endpoint.
\end{enumerate}
Steps 1--3 are proven within this paper (Step~3 via an explicit
closed-form formula for $H''(0)$ that is manifestly negative for all
$N \ge 3$).  Step 4 is established statistical mechanics.  Step 5 is a
physical observation.

The framework is universal: it does not depend on the specific dynamical
mechanism that produces the polygon.  Saturn's hexagon arises from Rossby
wave stationarity (a jet meander); Jupiter's cyclone rings arise from
discrete vortex self-organization (a cluster of storms).  Both mechanisms
deliver the system to $|\lam| \approx 1$, where the polygon lives---and
once there, eigenvalue branch disconnection and negative-temperature
statistics keep it.

The paper is organized as follows.  \Cref{sec:setup} introduces the
mathematical framework: the QGPV equation, M\"obius transformations, and
the loxodromic flow family.  \Cref{sec:theorems} proves the main results:
stationarity equivalence, branch disconnection, Thomson critical ratio,
amplitude formula, and energy--enstrophy sign structure.
\Cref{sec:verification} tests the framework against Saturn (Cassini) and
Jupiter (Juno) observations.  \Cref{sec:discussion} assembles the
five-step answer.
